% !TeX root = ../build/main.tex

We presented \davinci, a decentralized, verifiable voting protocol implemented as a specialized ZK rollup. By modeling elections as ZK-verified state transitions settled on Ethereum or on an Ethereum rollup, \davinci combines public auditability with scalable off-chain execution while minimizing reliance on any single party. Ballot secrecy is achieved through threshold ElGamal encryption under a key generated by a permissionless \ac{DKG} network or an invited council, while ciphertext re-randomization and silent revoting make revotes hard to tell from routine refreshes, which mitigates coercion and vote-selling.

The central change of this version of the protocol is the use of a \ac{zkVM}: the work of a sequencer, from the verification of every ballot proof, signature, and census proof to the update of the state, the re-encryption of the ballots, and the layout of the published data, is a single program whose execution is proven and verified on-chain with one constant-size proof. Expressing the transition function as a program made it practical to prove properties that the earlier circuit-based design left outside the proof, such as the exact content of the data blobs and the silent refresh of stored ballots, and it makes the statement that is enforced auditable as ordinary code. The data availability layer is decoupled from the settlement layer, so elections can settle on a low-cost rollup while their data remains on Ethereum.

To our knowledge, \davinci is the first system to combine a parametric ballot abstraction with a common additive tally, an open sequencer set, and fully proof-enforced election execution within a single rollup architecture. This enables a reusable voting infrastructure in which elections with different ballot formats and organizers can share the same settlement layer, verifier contracts, and proving pipeline.

Several directions remain open for future work. An important long-term direction is post-quantum migration, in particular the direct verification of \acp{STARK} and post-quantum signature and encryption schemes. Others are the chained proving mode, which folds a whole election into one on-chain verification, for elections that do not need per-batch settlement; the resharing of election keys across \ac{DKG} committees, so that a key outlives the committee that generated it; decentralized archival of blobs beyond their retention period; and a cast-as-intended mechanism for voters. We also plan to integrate with real-world credential providers and conduct large-scale deployment experiments in realistic electoral settings.
